% ============================================================
% 4 方法
% ============================================================
\section{方法}
\label{sec:method}

\subsection{设计动机：任务性质如何约束设计}
\label{sec:design}

本节先不描述结构，而是把第 3 节的任务性质逐条翻译成设计约束，
再在第 4.2--4.6 节给出对应实现；表~\ref{tab:design} 汇总这一对应关系，
并在右列标出每条设计在本章的验证位置。

\chead{（P1）来源信号在背景，不在人脸。}主流换脸算法保留目标视频的背景、头姿与动作，
只替换前景人脸；因此伪造帧与其源真实帧在背景上近乎逐像素相同。
第 5.9 节的 real$\rightarrow$real 对照进一步显示，换脸本身几乎不损失可溯源性
（$98.82\%$ 对 $98.36\%$），难度来自“同视频不同帧”与退化，而非“脸被换掉”。
据此，检索只应依赖\textbf{背景}这一在换脸后仍被保留的成分，
而不应依赖随篡改方法变化的人脸纹理或伪造痕迹——后者也与“无攻击监督”的设定冲突。

\chead{（P2）判别信号在中低分辨率的空间布局。}零训练扫描（第 5.2、5.3 节）显示：
不做任何训练、只保留 $8\times8$ 网格并逐图标准化的描述子，在帧级已达 $60.05\%$，
强于所有学习到的全局向量；继续提高网格分辨率（$16\times16$、$32\times32$）或只保留高频，
读数反而下降。这说明该任务需要的是\textbf{保留空间布局的局部描述子}，
把整图池化成一个全局向量会抹掉真正有用的那部分信息。

\chead{（P3）帧级难在“局部是否对得上”，而全局内积只回答“整体像不像”。}
同一条视频的相邻帧共享几乎全部背景，差异只出现在很小的人脸区域与轻微运动上。
两个全局向量的内积对这种差异不敏感——第 5.3 节把候选缩到同视频 $8$ 帧后，
学习编码器与零训练描述子的命中率持平（$59.79\%$ 对 $60.11\%$），
说明瓶颈不在描述子的容量，而在\textbf{比较方式}。
若两张图确实来自同一时刻，则它们的局部描述子在网格上应满足同一个几何变换；
于是把“整图比对”改为\textbf{成对局部匹配 $+$ 几何一致性打分}（第 4.4 节）。

\chead{（P4）几何打分的代价决定了它只能发生在候选层。}
逐对几何打分需要构造 $N\times N$ 的位置相关体（$256$px 输入下 $N=64^2=4096$），
全库两两计算在时间与显存上都不可行；而全局初筛在多数档位已把真值召回进前 $20$ 名。
因此采用\textbf{两阶段}结构：全库粗筛取 Top-$K$，只在候选内做几何重排（第 4.5 节）。
这也意味着系统能力受两个可分离的环节约束——初筛召回与候选内排序，
第 5.4 节用逐档诊断把二者分开测量。

\chead{（P5）两类退化的瓶颈不同，需要不同手段。}逐档诊断（第 5.4 节）给出：
\code{crop0.7off} 的损失 $100\%$ 来自初筛漏召回，
而 \code{jpeg5} 的初筛召回仍有 $99.3\%$、损失全在候选内打分（描述子在块效应下分数塌缩）。
前者需要一种对位移$/$缩放不敏感且廉价的\textbf{几何召回}，
后者需要一条在分数整体塌缩时\textbf{不被另一路拖累}的融合规则。
第 4.5 节据此给出几何回退与双路门控。

\chead{（P6）取证现场的源头多数不在库中。}因此检索之外还需\textbf{开集拒识}，
且判据只能使用无需标签的信号（对全库的 top-1 相似度，第 4.6 节）。

\begin{table}[t]
\centering
\caption{从任务性质到设计选择的对应关系。}
\label{tab:design}
\small
\begin{tabular}{p{3.2cm}p{4.6cm}p{3.2cm}c}
\toprule
任务性质 & 设计选择 & 针对的失效 & 验证 \\
\midrule
P1 背景被保留      & 以背景为检索对象，不建模人脸痕迹     & 人脸已被替换   & 5.2 / 5.9 \\
P2 信号在低频布局  & 保留 $64\times64$ 网格的局部描述子   & 全局池化抹布局 & 5.3 \\
P3 难点在局部对应  & 软互最近邻 $+$ 可微 DLT $+$ 软内点   & 全局内积不可校验 & 5.6 / 5.7 \\
P4 几何代价高      & 两阶段：全库 Top-$K$ $\rightarrow$ 候选内重排 & 全库几何不可行 & 5.6 \\
P5 两类退化瓶颈不同 & 双路融合 $+$ 门控；粗网格几何回退 & 分数塌缩 / 初筛漏召回 & 5.4 / 5.7 \\
P6 源头常不在库中  & top-1 相似度拒识                    & 无源查询       & 5.5 \\
\bottomrule
\end{tabular}
\end{table}

\subsection{整体框架}

系统由三部分组成：(i) 一个全局编码器，把整图映射为 $\ell_2$ 归一化的紧凑描述子，
只用于\textbf{全库初筛}（对应 P4）；(ii) 一对\textbf{局部描述子}，
在 $64\times64$ 网格上编码空间布局（对应 P2），供候选内的可微几何打分使用（对应 P3）；
(iii) 一个\textbf{开集拒识}判据（对应 P6）。
检索流程为：查询经全局编码器得到描述子 $\rightarrow$ 全库余弦初筛取 Top-$K$ 候选
$\rightarrow$ 候选内以双路局部描述子做几何重排，必要时触发门控与几何回退（对应 P5）
$\rightarrow$ 以 top-1 相似度判据决定接受或拒绝。

\subsection{全局编码器}

全局编码器采用多分支多尺度卷积主干，配合注意力特征融合（AFF）与嵌入头，
将整图映射为 $\ell_2$ 归一化描述子。训练采用监督对比损失
（SupCon\cite{khosla2020supcon}）与 VICReg\cite{bardes2022vicreg} 正则，
监督信号为源视频编号（粗粒度）或具体源帧编号（细粒度）。
该结构沿用已有工作的主干设计，本文不在结构层面提出新主张；
它在本文中的角色是提供一个可训练的初筛器，其读数在第 5.2、5.3 节与零训练描述子逐项对照。

\subsection{局部网格描述子与可微几何打分}
\label{sec:geo-module}

\textbf{局部描述子。}两个独立训练的卷积网络把图像编码为 $64\times64$ 的单位描述子场：
第一路为自建的小型主干（$64$ 维、宽 $64$），第二路取 ResNet-18 的浅层（$32$ 维、宽 $32$）。
两路都在同一网格上工作，训练目标是“同一空间位置对应同一物理点”，
因而可以直接比较，而无需检测关键点。选择网格而非稀疏关键点，是为了对应 P2：
栅格把空间布局完整保留下来，稀疏检测器在模糊与低分辨率下会先丢掉角点
（第 5.9、5.10 节观察到该档 SIFT 的中位好匹配数降为 $0$）。

\textbf{可微几何打分。}给定查询与候选的局部描述子场，$\mathrm{pair\_score}$ 由三步构成，
三步分别对应“找到对应—估计变换—衡量一致”这一几何验证链：

\begin{enumerate}[leftmargin=1.6em,itemsep=1pt]
\item \textbf{软互最近邻}。对两图的位置格计算相似度矩阵，以温度 $\tau$ 做行$/$列 softmax，
      得到双向的软匹配权重。用软匹配而不是硬 argmax，是为了让整条链路可导，
      同时避免 argmax 在小位移下把“几乎相同的两个位置”判成完全不同的匹配。
\item \textbf{可微 DLT}。以软匹配权重为权，求解单应矩阵的直接线性变换（DLT）最小二乘解。
      这一步把“局部对应是否服从同一个几何变换”显式写进打分，
      从而能区分“都在同一段背景里”与“确实是同一时刻的同一段背景”。
\item \textbf{软内点占比}。将候选点经单应变换后与原点的残差经尺度 $\sigma$ 映射为
      $[0,1]$ 的软内点权重，取其占比作为一致性度量。
\end{enumerate}

因此 $\mathrm{pair\_score}$ 对局部描述子与匹配参数均可导，可作为检索期打分器。
这条路径与经典 SIFT$+$RANSAC 构成两种可选的几何验证方案：前者不依赖特征检测器与随机采样，
后者则成熟且稳定。本文不预设何者更优，而是在同一几何退化协议、同一初筛下将二者并列比较
（第 5.6、5.7、5.9 节）。需要说明，这里的“可导”仅指几何打分内部
（软匹配 $\rightarrow$ DLT $\rightarrow$ 内点占比）；两路局部描述子当前为独立训练，
其梯度不回传全局编码器，把几何一致性并入端到端训练的尝试见第 5.8 节。

\subsection{两阶段重排器：融合、门控与几何回退}
\label{sec:two-stage}

\textbf{两阶段。}第一阶段以全局编码器做全库余弦初筛，取 Top-$K=20$ 候选；
第二阶段只在候选内计算 $\mathrm{pair\_score}$ 并重排。
这一分工使昂贵的几何计算与库规模解耦（P4），同时把误差分解为两个可独立测量的量：
初筛召回 $\mathrm{recall@}K$ 与候选内条件命中率。

\textbf{双路融合与门控。}两路描述子的失败模式互补：窄而深的第二路在模糊下更稳，
第一路在干净图上更锐。默认把两路分数逐查询 z-score 后以 $0.5:0.5$ 相加。
但极端压缩会让第一路的候选内分数整体塌缩成噪声，此时加权融合反而把噪声带进来；
于是引入一个\textbf{无标签}的门控：若第一路在候选内的分数标准差低于 $0.10$，
说明它已无力区分候选，则该查询完全改用第二路（P5 的“压缩”一侧）。

\textbf{几何回退。}非中心强裁剪会整体平移并缩放背景，使全局嵌入连“是哪个视频”都认不出
（第 5.4 节：\code{crop0.7off} 的视频级 recall@1 仅 $40.8\%$）。
这类失效无法靠候选内重排挽救，必须在召回阶段换一种相似度。
做法是：若一次重排后的最大几何分落在 $[0.35,0.68)$ 区间——即“最像的候选也不够像”，
说明真值大概没被召回——则把局部描述子池化到 $16\times16$ 粗网格，
用同一套可微几何打分对\textbf{全库}做一次廉价召回，取几何前 $10$ 与全局前 $10$ 的并集，
再做一次细网格重排。上下界缺一不可：分数低于 $0.35$ 说明描述子已整体塌缩
（如 \code{jpeg5}），此时低分与召回无关，回退反而有害（P5）。

上述三个部件都不引入新的训练，也不使用任何标签：两阶段、门控与回退的触发条件
全部是逐查询可观测的分数统计量。

\subsection{开集拒识}

真实取证中，待查伪造图的源头多数\textbf{不在库中}（P6）。
为此我们在检索前加入拒识判定：以查询对全库的 \textbf{top-1 余弦相似度}作为“库内成员度”打分，
低于阈值则拒绝。阈值在\textbf{开集样本分布}的目标 FPR 分位上标定（模拟以已知的“未知样本”确定工作点），
评测指标为拒识 AUC 与 TPR@FPR。
